\documentclass[a4paper,12pt]{article}
\usepackage[utf8]{inputenc}
\usepackage[T1]{fontenc}
\usepackage[ngerman]{babel}
\usepackage{amsmath}
\usepackage{amssymb}
\usepackage{enumitem}
\usepackage{array}
\usepackage{geometry}
\geometry{a4paper, margin=2.5cm}
\usepackage{fancyhdr}
\usepackage{parskip}

% Punkteangabe rechtsbündig am Ende eines Aufgabenteils
\newcommand{\pkt}[1]{\hfill\textit{(#1\,P)}}
% Kopf einer Aufgabe: \aufgabe{Nummer}{Titel}{Gesamtpunkte}
\newcommand{\aufgabe}[3]{%
  \bigskip\par\noindent\textbf{Aufgabe #1: #2}\hfill\textbf{#3 Punkte}%
  \par\nobreak\medskip}

\pagestyle{fancy}
\fancyhf{}
\lhead{\small 61211 Analysis -- Probeklausur 2}
\rhead{\small Matr.-Nr.: \makebox[3cm]{\dotfill}}
\cfoot{\small Seite \thepage}
\renewcommand{\headrulewidth}{0.4pt}

\begin{document}
\thispagestyle{empty}

\begin{center}
  {\LARGE\bfseries Probeklausur 2}\\[0.4em]
  {\large 61211 Analysis}\\[0.6em]
  {\normalsize FernUniversität in Hagen $\cdot$ Fakultät für Mathematik und Informatik}
\end{center}

\vspace{0.8em}
\noindent\small\textit{Diese Probeklausur ist bewusst beweisorientiert gehalten
(Normen, Funktionenfolgen, Differenzierbarkeit, \glqq beweise oder widerlege\grqq),
angelehnt an den Charakter der realen Klausuren.}\normalsize

\vspace{1.0em}

\noindent
\begin{tabular}{@{}ll@{}}
  Name: & \makebox[7cm]{\dotfill}\\[1.0em]
  Vorname: & \makebox[7cm]{\dotfill}\\[1.0em]
  Matrikelnummer: & \makebox[7cm]{\dotfill}\\
\end{tabular}

\vspace{1.2em}

\begin{center}
  \renewcommand{\arraystretch}{1.5}
  \begin{tabular}{|l|c|c|c|c|c|c|c||c|}
    \hline
    \textbf{Aufgabe} & 1 & 2 & 3 & 4 & 5 & 6 & 7 & \textbf{$\Sigma$}\\
    \hline
    max.\ Punkte & 20 & 16 & 14 & 14 & 12 & 10 & 14 & 100\\
    \hline
    erreicht & & & & & & & & \\
    \hline
  \end{tabular}
\end{center}

\vspace{0.8em}

\noindent\textbf{Bearbeitungshinweise}
\begin{itemize}[leftmargin=1.5em, itemsep=2pt]
  \item Die Bearbeitungszeit beträgt \textbf{120 Minuten}. Es sind alle 7 Aufgaben zu bearbeiten.
  \item Zugelassenes Hilfsmittel: ein handbeschriebenes DIN-A4-Blatt (auch beidseitig) mit eigenen Notizen.
  \item Verwenden Sie keinen Bleistift und keinen Rotstift.
  \item Formulieren Sie Ihre Beweise in vollständigen Sätzen und begründen Sie jeden Schritt. Verweise auf Sätze des Kurstextes dürfen benutzt werden.
\end{itemize}

\newpage

% =====================================================================
\aufgabe{1}{Multiple Choice}{20}

Kreuzen Sie an, ob die jeweilige Aussage \emph{wahr} oder \emph{falsch} ist. Für jede
richtig beantwortete Aussage erhalten Sie $2$ Punkte, für jede unbeantwortete Aussage
$1$ Punkt und für jede falsch beantwortete Aussage $0$ Punkte.

\medskip
\renewcommand{\arraystretch}{1.5}
\noindent\begin{tabular}{@{}c@{\hspace{0.6em}}p{11.4cm}@{\hspace{0.8em}}c@{\hspace{1.0em}}c@{}}
  & & \small Wahr & \small Falsch\\[0.2em]
  (a) & Jede beschränkte Folge in einem normierten Raum besitzt eine in diesem Raum konvergente Teilfolge. & $\square$ & $\square$\\
  (b) & In einem endlichdimensionalen normierten Raum konvergiert jede Cauchyfolge. & $\square$ & $\square$\\
  (c) & Jede stetige Funktion $f:[a,b]\to\mathbb{R}$ (mit $a<b$) ist gleichmäßig stetig. & $\square$ & $\square$\\
  (d) & Der punktweise Grenzwert einer Folge stetiger Funktionen $f_k:[0,1]\to\mathbb{R}$ ist stets stetig. & $\square$ & $\square$\\
  (e) & Existieren in $a\in\mathbb{R}^2$ die Richtungsableitungen $D_v f(a)$ für \emph{jeden} Einheitsvektor $v$, so ist $f$ in $a$ stetig. & $\square$ & $\square$\\
  (f) & Ist $f$ auf einer offenen Menge $U\subseteq\mathbb{R}^n$ stetig partiell differenzierbar, so ist $f$ auf $U$ total differenzierbar. & $\square$ & $\square$\\
  (g) & Gilt $\operatorname{grad} f(a)=0$ und ist die Hesse-Matrix $H(a)$ positiv semidefinit, so besitzt $f$ in $a$ ein lokales Minimum. & $\square$ & $\square$\\
  (h) & Besitzt $f$ unter der Nebenbedingung $g=0$ in $a$ ein lokales Extremum und ist $\operatorname{grad} g(a)\neq 0$, so gibt es ein $\lambda\in\mathbb{R}$ mit $\operatorname{grad} f(a)=\lambda\operatorname{grad} g(a)$. & $\square$ & $\square$\\
  (i) & Jede stetige Funktion $f:[0,\infty[\to\mathbb{R}$ mit $\lim_{x\to\infty}f(x)=0$ ist uneigentlich Riemann-integrierbar über $[0,\infty[$. & $\square$ & $\square$\\
  (j) & Jede stetig differenzierbare Kurve $\varphi:[a,b]\to\mathbb{R}^n$ ist rektifizierbar. & $\square$ & $\square$\\
\end{tabular}

% =====================================================================
\aufgabe{2}{Wann ist $N_p$ eine Norm?}{16}

Sei $n\ge 2$ und $p\in(0,1)$ fest. Betrachtet wird die Abbildung
\[
N_p:\mathbb{R}^n\to\mathbb{R},\qquad
N_p(x)=\Big(\sum_{k=1}^n |x_k|^p\Big)^{1/p},\qquad x={}^t(x_1,\dots,x_n).
\]
(Für $p\ge 1$ liefert dieselbe Formel nach Satz und Definition 1.4.4 eine Norm $\|\cdot\|_p$ auf $\mathbb{R}^n$; hier ist dagegen $0<p<1$.)
\begin{enumerate}[label=\alph*)]
  \item Zeigen Sie, dass $N_p$ die \emph{Definitheit} erfüllt, d.\,h.\ $N_p(x)\ge 0$ für alle $x$ und $N_p(x)=0\iff x=0$. \pkt{3}
  \item Zeigen Sie, dass $N_p$ \emph{positiv homogen} ist: $N_p(\alpha x)=|\alpha|\,N_p(x)$ für alle $\alpha\in\mathbb{R}$, $x\in\mathbb{R}^n$. \pkt{3}
  \item Weisen Sie anhand des Beispiels $x=e_1$, $y=e_2$ (Standardbasisvektoren) nach, dass die \emph{Dreiecksungleichung} für $N_p$ verletzt ist. Folgern Sie, dass $N_p$ für $p\in(0,1)$ \textbf{keine} Norm auf $\mathbb{R}^n$ ist. Wo geht im Vergleich zu $p\ge 1$ (Minkowski-Ungleichung, Satz 1.1.8) die Bedingung $p\ge 1$ ein? \pkt{4}
  \item Betrachten Sie nun $\displaystyle d(x,y):=\sum_{k=1}^n |x_k-y_k|^p$ (\emph{ohne} äußere $\tfrac1p$-Wurzel).
  \begin{enumerate}[label=(\roman*)]
    \item Beweisen Sie $(a+b)^p\le a^p+b^p$ für alle $a,b\ge 0$.
    \item Folgern Sie, dass $d$ eine \emph{Metrik} auf $\mathbb{R}^n$ ist (Definitheit, Symmetrie, Dreiecksungleichung, vgl.\ Definition und Satz 1.4.2).
    \item Begründen Sie kurz, dass $d$ dennoch von \emph{keiner} Norm erzeugt sein kann.
  \end{enumerate}
  \pkt{6}
\end{enumerate}

% =====================================================================
\aufgabe{3}{Funktionenfolge und gleichmäßige Konvergenz}{14}

Für $k\in\mathbb{N}$ sei
\[f_k:\mathbb{R}\to\mathbb{R},\qquad f_k(x):=\frac{kx}{1+k^2x^2}.\]
\begin{enumerate}[label=\alph*)]
  \item Bestimmen Sie die punktweise Grenzfunktion $f$ von $(f_k)$ auf $\mathbb{R}$. \pkt{3}
  \item Berechnen Sie $\|f_k-f\|_\infty=\sup_{x\in\mathbb{R}}|f_k(x)-f(x)|$ \emph{exakt} und entscheiden Sie, ob $(f_k)$ auf $\mathbb{R}$ gleichmäßig konvergiert. \pkt{5}
  \item Sei $a>0$ fest. Zeigen Sie, dass $(f_k)$ auf $[a,\infty[$ gleichmäßig gegen $f$ konvergiert. \pkt{4}
  \item Beweisen oder widerlegen Sie: \glqq Da alle $f_k$ stetig sind und die Grenzfunktion $f$ stetig ist, folgt aus dem Satz von Weierstraß 2.6.5, dass $(f_k)$ auf $\mathbb{R}$ gleichmäßig konvergiert.\grqq \pkt{2}
\end{enumerate}

% =====================================================================
\aufgabe{4}{Richtungsableitungen und Differenzierbarkeit}{14}

Sei $f:\mathbb{R}^2\to\mathbb{R}$ definiert durch
\[
f(x,y):=\begin{cases}\dfrac{x^3}{x^2+y^2},&(x,y)\neq(0,0),\\[1ex]0,&(x,y)=(0,0).\end{cases}
\]
\begin{enumerate}[label=\alph*)]
  \item Zeigen Sie, dass $f$ in $(0,0)$ stetig ist. \pkt{4}
  \item Zeigen Sie, dass für jeden Richtungsvektor $v=(v_1,v_2)^{T}$ mit $\|v\|_2=1$ die Richtungsableitung $D_v f(0,0)$ existiert, und berechnen Sie sie. Geben Sie insbesondere $\operatorname{grad} f(0,0)$ an. \pkt{6}
  \item Entscheiden Sie mit Begründung, ob $f$ in $(0,0)$ (total) differenzierbar ist. \pkt{4}
\end{enumerate}

% =====================================================================
\aufgabe{5}{Extremum unter einer Nebenbedingung}{12}

Für welches Paar positiver Zahlen $(x,y)$ mit $x+y=5$ wird $f(x,y)=x^3y^2$ maximal?
\begin{enumerate}[label=\alph*)]
  \item Bestimmen Sie mit der Methode der Lagrangemultiplikatoren (Satz 4.5.9) alle Punkte der Menge $M(g)=\{(x,y): x>0,\ y>0,\ x+y-5=0\}$, die als Extremalstelle von $f|_{M(g)}$ in Frage kommen, und geben Sie den zugehörigen Multiplikator $\lambda$ an. \pkt{8}
  \item Begründen Sie, dass $f|_{M(g)}$ dort tatsächlich ein globales Maximum annimmt. \emph{Hinweis:} Verhalten von $f$ am \glqq Rand\grqq{} $x\to0^+$ bzw. $y\to0^+$. \pkt{4}
\end{enumerate}

% =====================================================================
\aufgabe{6}{Beweise oder widerlege: Grenzwert und Integral}{10}

Beweisen oder widerlegen Sie die folgende Aussage:
\begin{quote}
\itshape Sei $(f_k)$ eine Folge in $R([0,1])$, die punktweise gegen $f\in R([0,1])$ konvergiert. Dann gilt stets $\displaystyle\lim_{k\to\infty}\int_0^1 f_k(x)\,dx=\int_0^1 f(x)\,dx$.
\end{quote}
Betrachten Sie zur Untersuchung die Folge $f_k:[0,1]\to\mathbb{R}$, $f_k(x):=k\,x\,(1-x^2)^k$.

% =====================================================================
\aufgabe{7}{Beweise oder widerlege: Kurvenintegrale}{14}

Beweisen oder widerlegen Sie:
\begin{enumerate}[label=\alph*)]
  \item Für jede stückweise glatte geschlossene Kurve $W$ in $\mathbb{R}^2$ und jedes auf einem Gebiet $G\supseteq|W|$ stetige Vektorfeld $f$ gilt $\displaystyle\int_W f\cdot dx=0$. \pkt{7}
  \item Erfüllt ein stetig differenzierbares Vektorfeld $f:G\to\mathbb{R}^n$ auf einem Gebiet $G$ die Integrabilitätsbedingungen $D_i f_j=D_j f_i$, so besitzt $f$ eine Stammfunktion. \pkt{7}
\end{enumerate}

\vfill
\begin{center}\small Viel Erfolg!\end{center}

\end{document}
